\section{The model and the frozen statement}\label{sec:model}

This section fixes the machine model and the statement~(T) exactly as they are formalized in the
frozen file \leanfile{Transpose/Statement.lean}, and introduces the vocabulary used in the rest of
the paper: total runs, obliviousness, correctness on a set of inputs, and co-entropy.

\subsection{Machines}

A machine $M$ has a finite state set $Q=\{0,\dots,q\}$ with initial state $0$, a read-only
two-way input tape, $k\ge 0$ read/write work tapes, and a write-only one-way output tape.
\begin{itemize}[leftmargin=2em]
\item The input tape holds the input word $w\in\{0,1\}^*$ in cells $0,\dots,\abs{w}-1$; every
other cell reads as an end marker.
\item Each work tape is two-way infinite over the alphabet $\Gamma=\{0,\dots,g\}$, where $0$ is the
blank.
\item In each step the machine may append at most one bit to the output.
\end{itemize}
The transition function $\delta$ maps the current state, the input symbol under the input head and
the $k$ work symbols under the work heads either to ``halt'' or to a new state, an input-head move,
a symbol to write and a head move for each work tape, and optionally one output bit. Head moves
are $-1$, $0$ or $+1$. All heads start at cell $0$, the work tapes start blank, and the output
starts empty. A configuration records the state, the head positions, the work-tape contents and the
output written so far. The machine \emph{computes $\mathit{out}$ on $w$ within $t$ steps}
(\lean{Transpose.ComputesWithin}) if, for some $s\le t$, the $s$-th configuration exists, is
halting, and has output exactly $\mathit{out}$.

Every step writes the cell under each work head (possibly rewriting the old symbol). The machine
is deterministic, and its transition table is finite; in particular $q$, $g$ and $k$ do not depend
on the input.

\subsection{Matrices and their encoding}

For $d\ge1$ (the side length of the matrix, written $n$ in the Lean files) let
\[
  b(d)=\max\bigl(1,\lceil\log_2 d\rceil\bigr)
  \qquad\text{(formally } \lfloor\log_2(d-1)\rfloor+1\text{ with }\lfloor\log_2 0\rfloor:=0\text{, \lean{Transpose.entryBits})}.
\]
A $d\times d$ matrix $A$ with $b$-bit entries is encoded row-major
(\lean{Transpose.encode}): the $d^2$ entries $(0,0),(0,1),\dots,(0,d-1),(1,0),\dots$ are written
as consecutive fixed-width blocks of $b$ bits, bit $0$ first, with no delimiters. The input length
is $L=d^2b$, and $d$ is determined by $L$. Bit $t$ of entry $(i,j)$ sits at position
$\alpha=(di+j)b+t$, and in the encoding of the transpose $A^{\T}$ the same bit sits at position
\[
  \operatorname{tpos}(\alpha)=(dj+i)b+t .
\]

\subsection{The statement}

\begin{definition}[Statement (T)]\label{def:T}
\formal{Statement.lean}{Transpose.TransposesWithin, Transpose.LinearTimeTransposable}
A machine $M$ \emph{transposes within $T(\cdot)$} if for every $d\ge1$ and every $d\times d$
matrix $A$ with $b(d)$-bit entries, $M$ computes the encoding of $A^{\T}$ on the encoding of $A$
within $T(L)$ steps, where $L=d^2b(d)$. Statement~(T), \lean{Transpose.LinearTimeTransposable},
says that some machine $M$ and some constant $c$ satisfy this with $T(L)=c\,(L+1)$.
\end{definition}

The file \leanfile{Transpose/Statement.lean} depends on Lean core only and is frozen: its SHA-256
hash, together with that of the checker \leanfile{Verify.lean}, is recorded in
\leanfile{FROZEN.sha256}. A resolution of~(T) is, by definition, a file
\leanfile{Transpose/Answer.lean} containing \texttt{theorem Transpose.answer} whose type is exactly
\lean{Transpose.LinearTimeTransposable} or its negation; \leanfile{Verify.lean} re-checks with the kernel, from scratch, every declaration of its environment
that does not ship with the Lean toolchain (\leanfile{Answer.lean} itself and the Mathlib and library modules it imports, directly or indirectly), and allows no axioms beyond
\texttt{propext}, \texttt{Classical.choice} and \texttt{Quot.sound}; in particular it rejects
\texttt{sorry}, \texttt{native\_decide} and new axioms. This paper does \emph{not} resolve~(T), so no \leanfile{Answer.lean} exists and
\leanfile{Verify.lean} does not compile yet; \cref{sec:conjecture} shows the exact form a proof of
the Tape Coding Conjecture would take as such a file.

\subsection{Runs}\label{sec:runs}

Throughout, the semantics of the frozen model is accessed through the \emph{total run}
\lean{run}: $\mathrm{run}(M,w,\tau)$ is the configuration after $\tau$ steps, frozen once the
machine halts. We write $\delta_\tau(w)$ for the transition taken at step $\tau$
(\lean{transAt}; it is ``halt'' exactly when the configuration at time $\tau$ is halting), and
say that $M$ \emph{halts at $T$ on $w$} (\lean{HaltsAt}) if steps $0,\dots,T-1$ execute and the
configuration at time $T$ is halting.

\begin{lemma}[The total run agrees with the frozen semantics]\label{lem:bridge}
\formal{Tape/Statements/Core.lean}{stepN_init_eq_some_iff, computesWithin_iff, haltsAt_unique, outRev_reverse_eq, transposes_of_transposesWithin}
\begin{enumerate}[label=(\alph*)]
\item The $\tau$-step iteration of the frozen step function from the initial configuration
succeeds with $c$ if and only if no step before $\tau$ halts and $\mathrm{run}(M,w,\tau)=c$.
\item $M$ computes $\mathit{out}$ on $w$ within $t$ steps if and only if $M$ halts on $w$ at some
$T\le t$ and the output of $\mathrm{run}(M,w,T)$ is $\mathit{out}$. The halting time is unique.
\item The output after $\tau$ steps is the list of the bits appended at steps $0,\dots,\tau-1$, in
order.
\item If $M$ transposes within any $T(\cdot)$, then for every $d\ge1$ it transposes $d\times d$
matrices with $b(d)$-bit entries in the sense of \cref{def:obl} below.
\end{enumerate}
\end{lemma}

\begin{outline}
By induction on the number of steps, comparing one step of the frozen semantics with one step of
the total run. Part~(d) is a restatement of the definitions.
\end{outline}

\begin{lemma}[Cell contents]\label{lem:cells}
\formal{Tape/Statements/Core.lean}{work_eq_written, work_eq_blank}
If work head $t$ is at cell $p$ at time $\sigma<\tau$ and not at $p$ at any time strictly between
$\sigma$ and $\tau$, then at time $\tau$ cell $p$ holds the symbol written at step $\sigma$ (after halting, the symbol
written is the one already in the cell, \lean{writtenSym}). A cell
not visited by its head before time $\tau$ is blank at time $\tau$.
\end{lemma}

\begin{outline}
Induction on $\tau$: a step changes only the cells under the work heads.
\end{outline}

\subsection{Obliviousness, correctness and co-entropy}\label{sec:obl}

\begin{definition}[Oblivious sets, transposers]\label{def:obl}
\formal{Tape/Defs.lean}{ObliviousOn, matrixInputs, Transposes}
Let $W$ be a set of inputs and $T\in\N$. The machine $M$ is \emph{oblivious on $W$ with halting
time $T$} if every $w\in W$ halts at $T$, and any two inputs of $W$ give the same input-head
position, the same work-head positions and the same decision whether to append an output bit, at
every time $\tau$ (including $\tau\ge T$). We write $\mathrm{Mat}(d,b)$ for the set of encodings of
$d\times d$ matrices with $b$-bit entries, and say that $M$ \emph{transposes $d\times d$ matrices
with $b$-bit entries} if on each of them it computes the encoding of the transpose within some
number of steps (no time bound).
\end{definition}

Since time $T$ is included, obliviousness also fixes the head moves of the last step; every
theorem below that assumes obliviousness assumes this (slightly stronger) form. By
\lean{matrixInputs_eq}, $\mathrm{Mat}(d,b)$ is exactly the set of words of length $d^2b$, so
``oblivious on $\mathrm{Mat}(d,b)$'' means oblivious at input length $L=d^2b$.

Many results hold for machines that are correct, or fast, only on a subset of the inputs. We
measure the size of such a subset by its co-entropy.

\paragraph{Co-entropy.}\label{def:coent} A set $W\subseteq\mathrm{Mat}(d,b)$ has \emph{co-entropy at most $m$} if
$2^{d^2b}\le 2^m\,\abs{W}$, that is, if $W$ contains at least a $2^{-m}$ fraction of all inputs. In the formal
statements this is written out as the hypothesis $2^{d^2b}\le2^m\cdot\abs{W}$ (named \texttt{hco}), with
$\abs{W}$ the cardinality \texttt{W.ncard}.

A machine is \emph{correct on $W$} (at time $T$) if every $w\in W$ encodes some matrix $A$ and the
output of $\mathrm{run}(M,w,T)$ is the encoding of $A^{\T}$. In the formal statements this
hypothesis is written out explicitly (it is the hypothesis \texttt{hcorr} of, for example,
\lean{cut_bound_on}).

\begin{lemma}[Encoding facts]\label{lem:encoding}
\formal{Tape/Statements/Core.lean}{matrixInputs_eq, tpos_spec}
The encodings of $d\times d$ matrices with $b$-bit entries are exactly the words of length $d^2b$.
For $j<d$ and $t<b$, $\operatorname{tpos}((di+j)b+t)=(dj+i)b+t$.
\end{lemma}

\begin{outline}
Induction on the nested concatenation defining the encoding, and integer division.
\end{outline}

\subsection{Non-vacuity}

The hypotheses of the lower bounds below (obliviousness together with correctness) are
satisfiable, and the counting quantities they involve are non-trivial on concrete machines.

\begin{proposition}[Witness machines]\label{prop:witnesses}
\formal{Tape/Statements/Core.lean}{copyTM_obliviousOn, copyTM_transposes, tr2_obliviousOn, tr2_transposes, tr2_cutPairs, tr2k_obliviousOn, tr2k_transposes}
\begin{enumerate}[label=(\alph*)]
\item The copy machine (which copies its input to its output) is oblivious on the inputs of every
length $\ell$, with halting time $\ell$, and transposes $1\times1$ matrices with $b$-bit entries
for every $b$.
\item A $9$-state machine \lean{Witness.tr2} without work tapes is oblivious on $\mathrm{Mat}(2,1)$
with halting time $5$ and transposes $2\times2$ matrices with $1$-bit entries. For the all-zero
input and the step set $J=\{2,3,4\}$ the set of cut pairs (\cref{def:cut}) is $\{(0,1)\}$.
\item The same machine with one idle work tape, \lean{Witness.tr2k} (so $k=1$), is oblivious on
$\mathrm{Mat}(2,1)$ with halting time $5$ and transposes $2\times2$ matrices with $1$-bit entries.
\end{enumerate}
\end{proposition}

\begin{outline}
For the copy machine, by induction on time, for inputs of every length $\ell$: at time $\tau\le\ell$
the input head is at cell $\tau$ and the output is the first $\tau$ input bits, and at time $\ell$ the
machine reads the end marker and halts. For \lean{Witness.tr2} and \lean{Witness.tr2k}, by evaluating
(with \texttt{decide}) the runs on all $16$ inputs of length $4$: the input-head positions $0,1,2,3,3,3$ (and, for \lean{Witness.tr2k}, the work-head position $0$) and the output schedule are common, the machine halts at time $5$, and the output on
$x_1x_2x_3x_4$ is $x_1x_3x_2x_4$. The cut pairs follow from the input-head positions $2,3,3$ at steps $2,3,4$ of the
all-zero run and the output positions $1,2,3$ appended at these steps.
\end{outline}
